\begin{examplebox}
	\textbf{\textcolor{CExample}{例 1（基础）}}

	\textbf{\textcolor{CExample}{题目：}}
	求函数 $y = 2\sin\left(3x + \frac{\pi}{4}\right)$ 的对称轴方程和对称中心。

	\textbf{\textcolor{CExample}{解题步骤：}}
	\begin{description}[style=nextline, leftmargin=2.8em, labelsep=0.5em, itemsep=0.45em]
		\item[\textbf{第一步（确定参数）：}]  $A=2,\ \omega=3,\ \varphi=\frac{\pi}{4}$。
			\par\textbf{理由：}明确函数参数以便套用公式。
		\item[\textbf{第二步（求对称轴）：}] 由对称轴条件 $3x + \frac{\pi}{4} = k\pi + \frac{\pi}{2}$，解得 $x = \frac{k\pi}{3} + \frac{\pi}{12}$（$k\in\mathbb{Z}$）。
			\par\textbf{理由：}直接代入对称轴公式 $x = \frac{k\pi + \pi/2 - \varphi}{\omega}$。
		\item[\textbf{第三步（求对称中心）：}] 由对称中心条件 $3x + \frac{\pi}{4} = k\pi$，解得 $x = \frac{k\pi}{3} - \frac{\pi}{12}$，故对称中心为 $\left(\frac{k\pi}{3} - \frac{\pi}{12}, 0\right)$（$k\in\mathbb{Z}$）。
			\par\textbf{理由：}直接代入对称中心公式 $x = \frac{k\pi - \varphi}{\omega}$，且纵坐标为 $0$。
	\end{description}

	\textbf{\textcolor{CExample}{关键结论：}}
	\textbf{对称轴为 $x = \frac{k\pi}{3} + \frac{\pi}{12}$，对称中心为 $\left(\frac{k\pi}{3} - \frac{\pi}{12}, 0\right)$（$k\in\mathbb{Z}$）。}

	\medskip
	\textbf{\textcolor{CExample}{例 2（稍变形）}}

	\textbf{\textcolor{CExample}{题目：}}
	求函数 $y = -3\sin\left(2x + \frac{\pi}{3}\right)$ 的对称轴方程和对称中心。

	\textbf{\textcolor{CExample}{解题步骤：}}
	\begin{description}[style=nextline, leftmargin=2.8em, labelsep=0.5em, itemsep=0.45em]
		\item[\textbf{第一步（确定参数）：}] $A=-3,\ \omega=2,\ \varphi=\frac{\pi}{3}$。注意 $A$ 的符号不影响对称性。
			\par\textbf{理由：}提取参数。
		\item[\textbf{第二步（求对称轴）：}] 由 $2x + \frac{\pi}{3} = k\pi + \frac{\pi}{2}$，解得 $x = \frac{k\pi}{2} + \frac{\pi}{12}$（$k\in\mathbb{Z}$）。
			\par\textbf{理由：}代入对称轴公式。
		\item[\textbf{第三步（求对称中心）：}] 由 $2x + \frac{\pi}{3} = k\pi$，解得 $x = \frac{k\pi}{2} - \frac{\pi}{6}$，对称中心为 $\left(\frac{k\pi}{2} - \frac{\pi}{6}, 0\right)$（$k\in\mathbb{Z}$）。
			\par\textbf{理由：}代入对称中心公式。
	\end{description}

	\textbf{\textcolor{CExample}{关键结论：}}
	\textbf{对称轴 $x = \frac{k\pi}{2} + \frac{\pi}{12}$，对称中心 $\left(\frac{k\pi}{2} - \frac{\pi}{6}, 0\right)$（$k\in\mathbb{Z}$）。}

	\medskip
	\textbf{\textcolor{CExample}{例 3（综合应用）}}

	\textbf{\textcolor{CExample}{题目：}}
	已知函数 $f(x) = 2\sin(2x + \varphi)$ 的一条对称轴为 $x = \frac{\pi}{3}$，求 $\varphi$ 的一个可能值。

	\textbf{\textcolor{CExample}{解题步骤：}}
	\begin{description}[style=nextline, leftmargin=2.8em, labelsep=0.5em, itemsep=0.45em]
		\item[\textbf{第一步（写出对称轴条件）：}] 对称轴满足 $2x + \varphi = k\pi + \frac{\pi}{2}$（$k\in\mathbb{Z}$）。
			\par\textbf{理由：}正弦型函数对称轴的相位条件。
		\item[\textbf{第二步（代入已知对称轴）：}] 将 $x = \frac{\pi}{3}$ 代入得 $2\cdot\frac{\pi}{3} + \varphi = k\pi + \frac{\pi}{2}$。
			\par\textbf{理由：}该点属于对称轴，满足方程。
		\item[\textbf{第三步（解出 $\varphi$）：}] 由 $\frac{2\pi}{3} + \varphi = k\pi + \frac{\pi}{2}$ 得
			\[
				\begin{aligned}
					\varphi &= k\pi + \frac{\pi}{2} \\
					& \quad - \frac{2\pi}{3} \\
					&= k\pi - \frac{\pi}{6}.
			\end{aligned}
			\]
			取 $k=0$ 得 $\varphi = -\frac{\pi}{6}$。
			\par\textbf{理由：}移项化简，$k\in\mathbb{Z}$。
	\end{description}

	\textbf{\textcolor{CExample}{关键结论：}}
	\textbf{$\varphi$ 的取值集合为 $\varphi = k\pi - \frac{\pi}{6}$（$k\in\mathbb{Z}$），一个可能值为 $-\frac{\pi}{6}$。}
\end{examplebox}
